%TODO: Ring, Order axioms and WOP; basic arithmetic and lemmas
%todo: ab \geq a for a, b in z+, and ab = a iff b is unit

\begin{definition}
Let $\bZ^+$ be the set of positive integers. We assume:
\begin{itemize}
    \item If $a,b\in\bZ^+$, then $a+b\in\bZ^+$.
    \item If $a,b\in\bZ^+$, then $ab\in\bZ^+$.
    \item $0\notin\bZ^+$.
    \item For every $a\in\bZ$, exactly one of $a\in\bZ^+$, $a=0$, or $-a\in\bZ^+$ is true.
\end{itemize}
\end{definition}

\begin{definition}
For $a,b\in\bZ$, define $a<b$ if $b-a\in\bZ^+$. Define $a>b$ if $b<a$. Define $a\leq b$ if $a<b$ or $a=b$, and define $a\geq b$ if $a>b$ or $a=b$.
\end{definition}

\begin{axiom}[Well-Ordering Principle]
Every nonempty subset of $\bZ^+$ has a least element.
\end{axiom}

\begin{lemma}
For $a,b,c\in\bZ$, if $a\leq b$ and $b\leq c$, then $a\leq c$.
\end{lemma}

\begin{proof}
If one of the inequalities is equality, the result follows directly. Now suppose $a<b$ and $b<c$. Then $b-a\in\bZ^+$ and $c-b\in\bZ^+$. By additive closure of $\bZ^+$,
\[
(b-a)+(c-b)\in\bZ^+.
\]
But $(b-a)+(c-b)=c-a$, so $c-a\in\bZ^+$. Thus $a<c$, and therefore $a\leq c$.
\end{proof}

\begin{lemma}[NIBZO]
There is no integer strictly between $0$ and $1$.
\end{lemma}

\begin{proof}
Suppose, for contradiction, that there exists $n\in\bZ$ such that $0<n<1$. Let
\[
S=\{n\in\bZ:0<n<1\}.
\]
Then $S$ is nonempty and $S\subseteq\bZ^+$. By WOP, $S$ has a least element. Call it $m$. Since $0<m$ and $m<1$, multiplying $m<1$ by the positive integer $m$ gives $m^2<m$. Also $m^2\in\bZ^+$ by multiplicative closure. Hence
\[
0<m^2<m<1.
\]
So $m^2\in S$, but $m^2<m$, contradicting the minimality of $m$. Therefore no such integer exists.
\end{proof}

\begin{lemma}
If $a\in\bZ^+$, then $1\leq a$.
\end{lemma}

\begin{proof}
Since $a\in\bZ^+$, we have $0<a$. If $a<1$, then $0<a<1$, contradicting NIBZO. Thus either $a=1$ or $1<a$, so $1\leq a$.
\end{proof}

\begin{lemma}
If $a,b\in\bZ^+$, then $ab\geq a$.
\end{lemma}

\begin{proof}
Since $b\in\bZ^+$, we have $1\leq b$. If $b=1$, then $ab=a$. If $1<b$, then $b-1\in\bZ^+$. Since $a\in\bZ^+$, multiplicative closure gives $a(b-1)\in\bZ^+$. But $a(b-1)=ab-a$, so $ab-a\in\bZ^+$, which means $a<ab$. Hence $a\leq ab$.
\end{proof}

\begin{lemma}
If $a,b\in\bZ^+$, then $ab=a$ if and only if $b=1$.
\end{lemma}

\begin{proof}
If $b=1$, then $ab=a$. Conversely, suppose $ab=a$. If $b\neq1$, then since $b\in\bZ^+$, the previous lemma gives $1<b$, so $b-1\in\bZ^+$. Since $a\in\bZ^+$, we get $a(b-1)\in\bZ^+$. But
\[
a(b-1)=ab-a=a-a=0,
\]
so $0\in\bZ^+$, contradicting the order axiom. Therefore $b=1$.
\end{proof}

\begin{lemma}[NIBZO]
There is no integer strictly between $0$ and $1$.  Equivalently, if
$n\in\mathbb{Z}$ and
\[
0<n<1,
\]
then such an $n$ cannot exist.
\end{lemma}

\begin{proof}
Suppose, for contradiction, that there is an integer strictly between $0$ and
$1$.  Put
\[
S=\{n\in\mathbb{Z}:0<n<1\}.
\]
Then $S$ is nonempty and $S\subseteq\mathbb{Z}_{\geq0}$.  By the well-ordering
principle, $S$ has a least element.  Call it $m$.

Since $0<m$, multiplicative closure of positive integers gives $0<m^2$.  Since
$m<1$ and $m>0$, multiplying the inequality $m<1$ by the positive integer $m$
gives
\[
m^2<m.
\]
Thus $m^2$ is also an integer satisfying
\[
0<m^2<m<1.
\]
So $m^2\in S$, but $m^2<m$.  This contradicts the choice of $m$ as the least
element of $S$.  Therefore $S$ is empty.
\end{proof}

\begin{definition}
    Define the \textbf{exponent} $\forall a \in \bZ, b \in \bZ \cup \{0\}$ as \[
    a^b = \begin{cases}
        1 & \text{if } b=0\\
        a \cdot a^{b-1} & \text{otherwise}
    \end{cases}
    \]
\end{definition}
For properties of exponents, see Appendix B.
